Let $G=(V,E)$ be a (simple, undirected) graph.

\begin{itemize}[leftmargin=*]
\item $\chi(G)$ denotes the chromatic number of $G$.
\item ``$G$ is triangle-free'' means $G$ contains no $K_3$ as a (not necessarily induced) subgraph.
\item For a graph $H$, we say $H$ is an \emph{induced subgraph} of $G$ if there is a vertex set $S\subseteq V$ such that $H\cong G[S]$, the subgraph induced by $S$.
\item Ambiguity: ``tree'' can mean finite or infinite. In the context of Gy\'{a}rf\'{a}s' induced-subtree conjecture, the intended meaning is \emph{finite tree}. We adopt:
\begin{quote}
A \emph{tree} $T$ is a finite, connected, acyclic graph.
\end{quote}
\end{itemize}

Question (interpreting ``tree'' as finite):
\begin{quote}
If $\chi(G)$ is infinite and $G$ is triangle-free, must it be that \emph{for every finite tree $T$} there exists an induced subgraph of $G$ isomorphic to $T$?
\end{quote}
